
%%%%%%%%%%%%%%%%%%%%%%%%%%%%%%%%%%%%%%%%%%%%%%%%%%%%%%%%%%%%%
\subsection{问题1：患者数据的标准化、清洗与管理}

\subsubsection{问题分析}

患者健康数据来源复杂、形式多样，存在以下主要挑战：数据来自不同医疗机构的不同信息系统（HIS、LIS、PACS等），存在缺失值、异常值、重复记录等数据质量问题；不同机构使用不同的数据编码和单位；日均数据量达万级规模。

本问采用数据工程方法进行求解，建立完整的数据ETL处理流程。

首先对多源患者健康数据进行采集和初步整理，接着对数据进行清洗处理（缺失值、异常值、重复记录），然后对数据进行标准化和转换（编码映射、格式统一、特征工程），最后将处理后的数据加载到统一的数据存储系统。

具体分析流程如图\ref{fig:analysis1_flow}所示。

\begin{figure}[H]
  \begin{center}
    \begin{tikzpicture}[x=1cm, y=0.7cm]
      \node[box] (n11) at (0, 0) {数据采集};
      \node[box] (n12) at (5, 0) {初步整理};
      \node[box] (n13) at (10, 0) {数据清洗};
      \node[box] (n22) at (5, -3) {标准化};
      \node[box] (n21) at (0, -3) {特征工程};
      \node[box] (n23) at (10, -3) {数据转换};
      \node[box] (n31) at (0, -6) {质量评估};
      \node[box] (n32) at (5, -6) {数据加载};
      \node[box] (n33) at (10, -6) {存储管理};
      
      \draw[arrow] (n11) -- (n12);
      \draw[arrow] (n12) -- (n13);
      \draw[arrow] (n13) -- (n23);
      \draw[arrow] (n23) -- (n22);
      \draw[arrow] (n22) -- (n21);
      \draw[arrow] (n21) -- (n31);
      \draw[arrow] (n31) -- (n32);
      \draw[arrow] (n32) -- (n33);
    \end{tikzpicture}
  \end{center}
  \caption{数据清洗与标准化处理流程}
  \label{fig:analysis1_flow}
\end{figure}

\subsubsection{模型建立}

针对数据质量问题，本节建立数据处理和质量评估模型。

\textbf{步骤1：缺失值处理}

对于不同类型的缺失数据采用不同的处理策略。对连续变量采用多元插补（MI）方法：
\begin{equation}
\hat{x}_{ij} = \bar{x_j} + \frac{\sum_{k=1}^{n}(x_{ik}-\bar{x_k})(x_{kj}-\bar{x_j})}{\sum_{k=1}^{n}(x_{ik}-\bar{x_k})^2} (x_{ij} - \bar{x_i})
\end{equation}

对分类变量采用众数或逻辑推理插补。对时间序列数据采用线性插值。

\textbf{步骤2：异常值检测与处理}

采用改进的基于孤立森林（Isolation Forest）的异常检测算法：
\begin{equation}
\text{异常得分} = 2^{-\frac{\text{平均路径长度}}{c(n)}}
\end{equation}

其中$c(n)$为样本数为$n$的平均路径长度：
\begin{equation}
c(n) = 2H(n-1) - \frac{2(n-1)}{n}
\end{equation}

其中$H(n)$为调和级数。异常得分>0.5的样本标记为异常。

\textbf{步骤3：数据标准化}

对数值变量进行Z-score标准化：
\begin{equation}
x'_{ij} = \frac{x_{ij} - \bar{x_j}}{\sigma_j}
\end{equation}

对分类变量进行one-hot编码。

\textbf{步骤4：数据质量评估}

定义综合质量指标：
\begin{equation}
Q_i = w_1 \cdot \text{完整率}_i + w_2 \cdot \text{准确率}_i + w_3 \cdot \text{一致性}_i
\end{equation}

其中权重$w_1=0.4, w_2=0.4, w_3=0.2$。

\subsubsection{求解结果}

数据处理结果如表\ref{tab:data_results}所示。系统能够处理日均万级数据量，处理速度达到10000条/秒以上，数据完整率和准确率均>99\%。
